\documentclass[twocolumn]{article}

\usepackage[margin=1in]{geometry}
\usepackage{amsmath,amssymb,amsthm}
\usepackage{graphicx}
\usepackage{booktabs}
\usepackage{hyperref}
\usepackage{xcolor}
\usepackage{algorithm}
\usepackage{algpseudocode}
\usepackage{caption}
\usepackage{subcaption}
\usepackage{multirow}
\usepackage{url}
\usepackage{cite}
\usepackage{braket}  % Dirac notation

\hypersetup{colorlinks=true, linkcolor=blue, citecolor=blue, urlcolor=blue}

\title{\textbf{EQ-GAN: Entanglement-Enhanced Quantum Generative Adversarial Network\\
with Conditional Style Control}}

\author{
  Felipe Mahlow\\
  \small \texttt{felipe.mahlow@gmail.com}
}

\date{\today}

\begin{document}

\maketitle

% ============================================================
\begin{abstract}
Quantum Generative Adversarial Networks (QGANs) have emerged as a promising
framework for quantum-enhanced generative modeling, yet existing architectures
treat parametric quantum circuits (PQCs) primarily as generic function
approximators, leaving quantum entanglement largely unexploited as a
computational resource. We propose \textbf{EQ-GAN}, an entanglement-enhanced
hybrid QGAN architecture that introduces \emph{quantum entanglement as an
explicit style injection mechanism} for conditional image generation. EQ-GAN
maintains two quantum registers: a content register (6 qubits, angle-encoded
from latent noise) and a style register (2 qubits, IQP feature-mapped from
class labels), connected by an entanglement core consisting of a structured
CNOT ladder that creates cross-register correlations before variational layers
are applied. The circuit output drives a classical convolutional decoder to
produce 28$\times$28 images. We evaluate EQ-GAN on FashionMNIST, MNIST, and
BreastMNIST using FID, Inception Score, and Precision\,/\,Recall metrics,
and benchmark against DCGAN, WGAN-GP, and prior QGAN architectures.
A systematic ablation study isolates the contribution of the entanglement
core by comparing against variants with independent registers and no style
register. We further characterize the circuit via expressibility analysis
(KL divergence from the Haar measure) and von~Neumann entanglement entropy
tracking during training. Our results provide quantitative evidence on
when and how quantum entanglement contributes to generative quality,
advancing the empirical understanding of entanglement as a computational
resource in near-term quantum machine learning.
\end{abstract}

\textbf{Keywords:} quantum GAN, entanglement, hybrid quantum-classical,
conditional image generation, NISQ, parametric quantum circuits

% ============================================================
\section{Introduction}
\label{sec:intro}

Generative Adversarial Networks (GANs) \cite{goodfellow2014generative} have
transformed the landscape of generative modeling, enabling the synthesis of
high-fidelity images through an adversarial training framework between a
generator $G$ and discriminator $D$. Classical advances—including the
Wasserstein formulation \cite{arjovsky2017wasserstein}, style-based
architectures \cite{karras2019style}, and self-attention mechanisms
\cite{zhang2019self}—have progressively improved image quality, diversity,
and controllability.

Simultaneously, the emergence of near-term quantum processors has motivated
the study of quantum-enhanced generative models. The theoretical appeal is
rooted in the exponential dimensionality of Hilbert space: an $n$-qubit
system spans a $2^n$-dimensional state space, potentially offering a
representational advantage for capturing complex distributions
\cite{biamonte2017quantum}. Quantum GANs (QGANs) were first proposed by
Lloyd and Weedbrook \cite{lloyd2018quantum} and Dallaire-Demers and Killoran
\cite{dallaire2018quantum}, establishing the theoretical foundations for
adversarial training with quantum generators. Subsequent hybrid architectures
have demonstrated practical generation of low-dimensional data on near-term
hardware \cite{huang2021experimental, tsang2023mosaiq, rudolph2024latent}.

However, a critical gap persists in existing QGAN architectures: \emph{quantum
entanglement is not exploited as an explicit mechanism for generative control}.
Most hybrid QGANs employ PQCs as generic function approximators, using
entanglement implicitly (through CNOT-rich ansatze) without differentiating
between the roles of content generation and style conditioning in quantum
space. This stands in contrast to classical advances such as StyleGAN
\cite{karras2019style}, where disentangled style codes enable precise control
over generative attributes through adaptive instance normalization.

We propose \textbf{EQ-GAN} (Entanglement-enhanced Quantum GAN), which
addresses this gap by introducing a dedicated \emph{entanglement core}—a
structured CNOT ladder that connects a style qubit register (encoding class
information via IQP feature maps) to a content qubit register (encoding
latent noise via angle encoding). The central hypothesis is that
quantum entanglement between these registers creates cross-register
correlations that cannot be replicated by independent classical conditioning,
enabling more effective class-conditional generation in the quantum regime.

\paragraph{Contributions.} This paper makes four contributions:
\begin{enumerate}
  \item \textbf{EQ-GAN architecture}: a hybrid QGAN with 8-qubit circuit
    (6 content, 2 style) featuring a structured entanglement core for
    conditional image generation.
  \item \textbf{Ablation study}: systematic isolation of the entanglement
    core's contribution by comparing full EQ-GAN against independent-register
    and unconditional variants.
  \item \textbf{Quantum characterization}: expressibility analysis and
    von~Neumann entanglement entropy tracking during training.
  \item \textbf{Empirical benchmark}: evaluation on FashionMNIST, MNIST, and
    BreastMNIST with FID, IS, and Precision/Recall metrics against classical
    and quantum baselines.
\end{enumerate}

The remainder of the paper is organized as follows. Section~\ref{sec:related}
reviews related work. Section~\ref{sec:method} describes the EQ-GAN
architecture. Section~\ref{sec:theory} presents the theoretical
characterization. Section~\ref{sec:experiments} details the experimental
setup. Section~\ref{sec:results} presents and discusses results.
Section~\ref{sec:conclusion} concludes with limitations and future directions.

% ============================================================
\section{Related Work}
\label{sec:related}

\subsection{Classical GANs and Style Control}

The original GAN \cite{goodfellow2014generative} frames generative modeling
as a minimax game: $\min_G \max_D \mathbb{E}[\log D(x)] +
\mathbb{E}[\log(1-D(G(z)))]$. Conditional GANs \cite{mirza2014conditional}
introduced class-conditional generation by concatenating label information to
both generator and discriminator inputs. The Wasserstein GAN with gradient
penalty (WGAN-GP) \cite{gulrajani2017improved} replaced the Jensen-Shannon
divergence with the Earth Mover's distance, providing stable training dynamics.

Style-based generation \cite{karras2019style, karras2020analyzing} introduced
a fundamentally different conditioning paradigm: instead of concatenating class
labels, a learned mapping network transforms latent codes into \emph{style
parameters} that are injected at each layer via adaptive instance normalization.
This disentanglement of style and content has been highly influential—our
entanglement core is conceptually motivated by this paradigm, transposing it
into quantum space via CNOT-based cross-register correlations.

\subsection{Quantum Generative Adversarial Networks}

Lloyd and Weedbrook \cite{lloyd2018quantum} first proposed a fully quantum GAN
where both generator and discriminator are quantum channels, establishing
theoretical bounds on training dynamics and convergence. Dallaire-Demers and
Killoran \cite{dallaire2018quantum} introduced a framework for training QGANs
on quantum hardware using a quantum cost function based on the diamond norm.

Hybrid architectures—combining quantum generators with classical
discriminators—have dominated practical QGAN research. Huang et al.\
\cite{huang2021experimental} demonstrated a patch-based hybrid QGAN on
superconducting hardware, achieving generation of $8\times 8$ handwritten
digits by assembling image patches from separate PQC outputs. The MosaIQ
architecture \cite{tsang2023mosaiq} extended this with a mosaic strategy,
generating overlapping patches and aggregating them to reduce boundary
artifacts.

Latent QGANs \cite{rudolph2024latent} operate in a compressed latent space
produced by a classical autoencoder, reducing the dimensionality challenge
while enabling the quantum circuit to focus on the semantic structure of the
data. Recent work has explored quantum Wasserstein distances
\cite{kiani2022quantum}, quantum normalizing flows
\cite{johansen2024quantum}, and barren plateau mitigation strategies
\cite{mcclean2018barren, zhang2024trainability} relevant to QGAN training.

Our work differs from existing hybrid QGANs in that we explicitly differentiate
the roles of content and style registers and use entanglement as a structured
conditioning mechanism rather than incidental circuit topology.

\subsection{Entanglement as a Computational Resource}

Entanglement entropy has been studied as a measure of quantum circuit
expressibility \cite{sim2019expressibility}. Abbas et al.\
\cite{abbas2021power} demonstrated that quantum models can achieve higher
effective dimension than comparable classical models. Bipartite entanglement
in PQCs has been linked to trainability \cite{cerezo2021cost} and generalization
\cite{caro2022generalization} in quantum machine learning.

In the context of generative models, entanglement enables the generation of
states with non-classical correlations \cite{lloyd2018quantum}. Our work
provides empirical evidence connecting entanglement entropy to generative
quality in a hybrid setting, operationalizing entanglement as a style
injection mechanism.

% ============================================================
\section{EQ-GAN Architecture}
\label{sec:method}

\subsection{Overview}

EQ-GAN is a conditional hybrid QGAN with five components: (1) a
classical style encoder mapping class labels to quantum style inputs,
(2) a classical noise encoder mapping latent noise to quantum content inputs,
(3) a quantum circuit implementing content encoding, IQP-based style encoding,
an entanglement core, and variational layers,
(4) a classical convolutional decoder mapping quantum measurements to images, and
(5) a conditional CNN discriminator. The full architecture is shown in
Figure~\ref{fig:architecture}.

\subsection{Classical Encoders}

\paragraph{Noise encoder.}
Given Gaussian noise $\mathbf{z} \sim \mathcal{N}(\mathbf{0}, \mathbf{I}_{64})$,
the noise encoder $E_\text{noise}: \mathbb{R}^{64} \rightarrow [-\pi, \pi]^{N_C}$
is a two-layer MLP that produces $N_C = 6$ angle parameters $\boldsymbol{\theta}_n$
scaled to $[-\pi, \pi]$ via $\tanh(\cdot) \times \pi$:
\begin{equation}
  \boldsymbol{\theta}_n = \pi \cdot \tanh(W_2 \cdot \text{LeakyReLU}(W_1 \mathbf{z} + b_1) + b_2).
\end{equation}

\paragraph{Style encoder.}
Given a one-hot class label $\mathbf{c} \in \{0,1\}^{N_\text{cls}}$, the style
encoder $E_\text{style}: \mathbb{R}^{N_\text{cls}} \rightarrow [-\pi,\pi]^{N_S}$
produces $N_S = 2$ style parameters $\boldsymbol{\phi}_s$ via a similar two-layer
MLP.

\subsection{Quantum Circuit}
\label{sec:circuit}

The quantum circuit operates on $N_Q = 8$ qubits partitioned into a
\emph{content register} $\mathcal{C} = \{q_0, \ldots, q_5\}$ and a
\emph{style register} $\mathcal{S} = \{q_6, q_7\}$.

\paragraph{Content encoding.} Angle encoding applies individual $R_Y$ rotations
to the content register:
\begin{equation}
  U_n(\boldsymbol{\theta}_n) = \bigotimes_{i=0}^{5} R_Y(\theta_{n,i}).
\end{equation}

\paragraph{Style encoding (IQP feature map).}
The style register uses an IQP-inspired feature map
\cite{havlicek2019supervised}, which is known to produce states that are
classically hard to simulate for random inputs:
\begin{equation}
  U_s(\boldsymbol{\phi}_s) =
    \text{IsingZZ}(\phi_{s,0} \cdot \phi_{s,1}, \{q_6, q_7\}) \cdot
    \bigotimes_{j=6}^{7} [R_Z(\phi_{s,j-6}) \cdot H_{q_j}].
\end{equation}

\paragraph{Entanglement core.}
The entanglement core $U_E$ is the key architectural contribution. It consists
of a CNOT ladder connecting the style register to the content register,
followed by intra-style entanglement:
\begin{equation}
  U_E = \text{CNOT}_{q_7 \to q_6} \cdot \text{CNOT}_{q_6 \to q_0} \cdot \text{CNOT}_{q_7 \to q_1},
  \label{eq:entcore}
\end{equation}
where CNOT$_{a \to b}$ denotes a CNOT gate with control qubit $a$ and target
qubit $b$. This structure creates quantum correlations between the style and
content registers, entangling the class information encoded in $\mathcal{S}$
with the content generation subspace $\mathcal{C}$.

\paragraph{Variational layers.}
A hardware-efficient ansatz \cite{kandala2017hardware} with $L = 3$ layers
is applied to all 8 qubits:
\begin{equation}
  U_V(\boldsymbol{\omega}) = \prod_{\ell=1}^{L}
    \left[ U_\text{CNOT\,ring} \cdot \bigotimes_{i=0}^{7}
    R_Z(\omega_{\ell,i,2}) R_Y(\omega_{\ell,i,1}) R_Z(\omega_{\ell,i,0}) \right],
\end{equation}
where $\boldsymbol{\omega} \in \mathbb{R}^{L \times N_Q \times 3}$ are
trainable variational parameters ($72$ parameters total).

The full circuit produces the state:
\begin{equation}
  \ket{\psi(\boldsymbol{\theta}_n, \boldsymbol{\phi}_s, \boldsymbol{\omega})} =
    U_V \cdot U_E \cdot [U_n \otimes U_s] \ket{0}^{\otimes 8},
\end{equation}
and the generator output is the vector of Pauli-$Z$ expectation values on the
content register:
\begin{equation}
  \mathbf{q} = (\bra{\psi} Z_0 \ket{\psi}, \ldots, \bra{\psi} Z_5 \ket{\psi})
    \in [-1, 1]^6.
\end{equation}

\subsection{Classical Decoder}
The decoder $D_\text{dec}: [-1,1]^6 \rightarrow \mathbb{R}^{1\times 28\times 28}$
maps quantum measurements to images via a fully-connected layer followed by
three transposed convolutional layers (4$\times$4→7$\times$7→14$\times$14→28$\times$28)
with batch normalization and ReLU activations, and a final $\tanh$ output.

\subsection{Discriminator and Training Objective}
The discriminator $D$ is a conditional CNN that concatenates a learned class
embedding (projected to a spatial map) with the input image, followed by three
strided convolutional layers with instance normalization.

EQ-GAN is trained with the WGAN-GP objective \cite{gulrajani2017improved}:
\begin{align}
  \mathcal{L}_D &= \mathbb{E}[D(G(\mathbf{z},\mathbf{c}))] - \mathbb{E}[D(\mathbf{x})]
    + \lambda \, \mathbb{E}\!\left[(\|\nabla D(\hat{\mathbf{x}})\|_2 - 1)^2\right], \\
  \mathcal{L}_G &= -\mathbb{E}[D(G(\mathbf{z},\mathbf{c}))],
\end{align}
where $\hat{\mathbf{x}}$ is a convex interpolation between real and fake samples,
and $\lambda = 10$ is the gradient penalty coefficient.

Gradients with respect to the variational parameters $\boldsymbol{\omega}$ are
computed via the adjoint differentiation method \cite{jones2020differentiating},
which is exact and parameter-efficient.

\subsection{Ablation Variants}

To isolate the entanglement core's contribution, we define two ablation models:

\begin{itemize}
  \item \textbf{EQ-GAN$_\text{no-ent}$}: identical to EQ-GAN but with $U_E$
    removed. Content and style registers evolve independently before the
    variational layers, with no cross-register quantum correlations.
  \item \textbf{EQ-GAN$_\text{no-style}$}: unconditional variant using only the
    6-qubit content register with no style input, analogous to a standard QGAN.
\end{itemize}

% ============================================================
\section{Theoretical Characterization}
\label{sec:theory}

\subsection{Expressibility}

Following Sim et al.\ \cite{sim2019expressibility}, we quantify circuit
expressibility as the KL divergence between the distribution of pairwise
state fidelities $|\langle\psi(\boldsymbol{\theta}_1)|\psi(\boldsymbol{\theta}_2)\rangle|^2$
for uniformly random parameters and the Haar-random distribution
$F(f) = (2^n - 1)(1-f)^{2^n - 2}$:
\begin{equation}
  \mathcal{E} = D_\text{KL}\!\left(P_\text{circuit}(f) \,\|\, P_\text{Haar}(f)\right).
\end{equation}
Lower $\mathcal{E}$ indicates higher expressibility. We compare
$\mathcal{E}_\text{EQ-GAN}$ against $\mathcal{E}_\text{no-ent}$ to assess
whether the entanglement core increases the expressibility of the circuit.

\subsection{Entanglement Entropy}

We track the von~Neumann entanglement entropy of the bipartition
$\mathcal{C} | \mathcal{S}$ during training:
\begin{equation}
  S(\rho_\mathcal{C}) = -\text{Tr}(\rho_\mathcal{C} \log_2 \rho_\mathcal{C}),
\end{equation}
where $\rho_\mathcal{C} = \text{Tr}_\mathcal{S}[\ket{\psi}\bra{\psi}]$ is the
reduced density matrix of the content register. If the entanglement core plays
an active role in generation, we expect $S(\rho_\mathcal{C})$ to increase
during training as the variational layers learn to leverage the entanglement
created by $U_E$.

% ============================================================
\section{Experiments}
\label{sec:experiments}

\subsection{Datasets}

We evaluate on three benchmark datasets:
\begin{itemize}
  \item \textbf{FashionMNIST} \cite{xiao2017fashion} (primary): 10 classes,
    $28\times28$ grayscale, 5{,}000 training samples used.
  \item \textbf{MNIST} \cite{lecun1998gradient} (sanity check): 10 classes,
    $28\times28$ grayscale, 5{,}000 training samples.
  \item \textbf{BreastMNIST} \cite{medmnistv2} (application): 2-class
    ultrasound (benign vs.\ malignant), $28\times28$ grayscale, from
    MedMNIST v2.
\end{itemize}
All images are normalized to $[-1, 1]$. Training is conducted on 5{,}000
samples to ensure comparability across quantum models (a standard subset size
in QGAN literature \cite{huang2021experimental, tsang2023mosaiq}).

\subsection{Baselines}

\paragraph{Classical baselines.}
DCGAN \cite{radford2015unsupervised} and WGAN-GP \cite{gulrajani2017improved}
provide classical references, both trained with identical data splits and
epoch counts.

\paragraph{Quantum baselines.}
We implement the following quantum baselines following their original
formulations:
(i) Patch QGAN \cite{huang2021experimental} — 4-qubit patch-based generator,
(ii) MosaIQ \cite{tsang2023mosaiq} — mosaic patch strategy,
(iii) Latent QGAN \cite{rudolph2024latent} — quantum generator in VAE latent space.

\paragraph{Ablations.}
EQ-GAN$_\text{no-ent}$ and EQ-GAN$_\text{no-style}$ as defined in
Section~\ref{sec:method}.

\subsection{Evaluation Metrics}

\begin{itemize}
  \item \textbf{FID} \cite{heusel2017gans}: Fréchet Inception Distance between
    real and generated InceptionV3 feature distributions (lower is better).
  \item \textbf{IS} \cite{salimans2016improved}: Inception Score,
    measuring diversity and quality (higher is better).
  \item \textbf{Precision / Recall} \cite{kynkaanniemi2019improved}: fraction
    of generated samples in the real manifold (precision) and fraction of
    real samples covered (recall).
  \item \textbf{Entanglement Entropy}: $S(\rho_\mathcal{C})$ averaged over
    50 random input states at each checkpoint.
  \item \textbf{Expressibility} $\mathcal{E}$: estimated from 1{,}000
    random parameter pairs.
  \item \textbf{Parameter count}: total trainable parameters.
\end{itemize}

FID and IS are evaluated on $N_\text{eval} = 2{,}000$ generated images
using InceptionV3 \cite{szegedy2016rethinking} features.

\subsection{Implementation Details}

EQ-GAN is implemented in PennyLane 0.45 \cite{bergmann2018pennylane} with
PyTorch 2.4 \cite{paszke2019pytorch}. The \texttt{lightning.qubit} backend
provides exact state-vector simulation. Quantum gradients are computed via
the adjoint method \cite{jones2020differentiating}. Classical components use
GPU acceleration (NVIDIA RTX 3080 Ti, CUDA 13.0); quantum simulation runs on
CPU. Training uses the Adam optimizer
($\beta_1 = 0, \beta_2 = 0.9$) for both generator and discriminator, with
learning rates $\eta_G = \eta_D = 10^{-4}$, batch size 32, $n_\text{critic} = 3$,
and $\lambda_\text{GP} = 10$, for 50 epochs.

% ============================================================
\section{Results and Discussion}
\label{sec:results}

\subsection{Quantitative Results}

Table~\ref{tab:main_results} reports FID, IS, Precision, and Recall for all
models on FashionMNIST, MNIST, and BreastMNIST. [Results will be populated
upon experiment completion.]

\begin{table*}[t]
  \centering
  \caption{Main results: FID ($\downarrow$), IS ($\uparrow$), Precision ($\uparrow$),
    and Recall ($\uparrow$) at epoch 50. Best QGAN in \textbf{bold}.
    $\dagger$: classical reference (not directly comparable).}
  \label{tab:main_results}
  \small
  \begin{tabular}{llcccccccc}
    \toprule
    & & \multicolumn{4}{c}{FashionMNIST} & \multicolumn{2}{c}{MNIST} & \multicolumn{2}{c}{BreastMNIST} \\
    \cmidrule(lr){3-6}\cmidrule(lr){7-8}\cmidrule(lr){9-10}
    Type & Model & FID$\downarrow$ & IS$\uparrow$ & P$\uparrow$ & R$\uparrow$
         & FID$\downarrow$ & IS$\uparrow$ & FID$\downarrow$ & IS$\uparrow$ \\
    \midrule
    \multirow{2}{*}{Classical$^\dagger$}
      & DCGAN   & -- & -- & -- & -- & -- & -- & -- & -- \\
      & WGAN-GP & -- & -- & -- & -- & -- & -- & -- & -- \\
    \midrule
    \multirow{3}{*}{Quantum}
      & Patch QGAN  & -- & -- & -- & -- & -- & -- & -- & -- \\
      & MosaIQ      & -- & -- & -- & -- & -- & -- & -- & -- \\
      & Latent QGAN & -- & -- & -- & -- & -- & -- & -- & -- \\
    \midrule
    \multirow{3}{*}{Ablation}
      & EQ-GAN$_\text{no-style}$ & -- & -- & -- & -- & -- & -- & -- & -- \\
      & EQ-GAN$_\text{no-ent}$   & -- & -- & -- & -- & -- & -- & -- & -- \\
      & \textbf{EQ-GAN (ours)}  & -- & -- & -- & -- & -- & -- & -- & -- \\
    \bottomrule
  \end{tabular}
\end{table*}

\subsection{Ablation Study}

Figure~\ref{fig:ablation} compares FID across ablation variants.
[To be populated.]

\subsection{Quantum Characterization}

\paragraph{Expressibility.}
Table~\ref{tab:expressibility} compares the KL divergence $\mathcal{E}$
for EQ-GAN and EQ-GAN$_\text{no-ent}$.

\paragraph{Entanglement entropy.}
Figure~\ref{fig:entropy} shows the evolution of $S(\rho_\mathcal{C})$ during
training, indicating whether the variational layers learn to leverage the
entanglement structure created by $U_E$.

\subsection{Qualitative Results}

Figure~\ref{fig:samples} shows generated samples for all 10 FashionMNIST
classes at epochs 10, 30, and 50. [To be populated.]

\subsection{Discussion}

\paragraph{Entanglement contribution.}
The ablation between EQ-GAN and EQ-GAN$_\text{no-ent}$ directly tests whether
the CNOT-based style injection is the source of any generative improvement.
If EQ-GAN outperforms EQ-GAN$_\text{no-ent}$ with statistical significance,
the entanglement core is a net contributor to generative quality. If not, the
improvement is attributable to the IQP-encoded style information alone, and
the entanglement provides no additional benefit.

\paragraph{NISQ limitations.}
All results are from classical simulation of quantum circuits. The findings
characterize the \emph{algorithmic} role of entanglement in the architecture,
independent of hardware noise. Running EQ-GAN on actual quantum processors
would introduce gate errors and decoherence, potentially degrading generation
quality. Noise-mitigation strategies \cite{temme2017error} could be explored
in future work.

\paragraph{Parameter efficiency.}
EQ-GAN contains only 72 variational quantum parameters out of
$\sim$2.8M total, with the remainder in the classical encoder-decoder.
This ratio reflects the dominant classical contribution and raises the question
of whether the quantum components provide genuine value or are overshadowed
by the classical decoder's capacity.

% ============================================================
\section{Conclusion}
\label{sec:conclusion}

We presented EQ-GAN, a conditional hybrid QGAN architecture that introduces
quantum entanglement as an explicit style injection mechanism via a structured
CNOT entanglement core connecting dedicated content and style qubit registers.
Through systematic evaluation on three image datasets and ablation studies
isolating the entanglement contribution, we provide quantitative empirical
evidence on the role of entanglement in near-term quantum generative models.

Our findings contribute to an ongoing empirical program of understanding
quantum advantages and limitations in the NISQ era, providing an architectural
blueprint for entanglement-aware QGAN design and a reproducible benchmark for
future quantum generative models.

\paragraph{Future work.} Promising directions include: (i) evaluation on
actual quantum processors to assess the impact of hardware noise; (ii) scaling
to higher-resolution images via hierarchical quantum-classical architectures;
(iii) extending the entanglement core to multi-class style control with
larger style registers; and (iv) theoretical analysis of the expressibility
advantage conferred by the entanglement core.

% ============================================================
\paragraph{Limitations and ethical considerations.}
This work does not demonstrate quantum advantage over classical methods.
All experiments were conducted on simulated quantum circuits. BreastMNIST
experiments are conducted as a generalization benchmark only; the generated
medical images should not be used for clinical purposes.

% ============================================================
\bibliographystyle{unsrt}
\bibliography{references}

\end{document}
